\clearpage
\fancyhead[L]{\small SOLDNER. AZIMUTHE ZU REDUCIREN.}
\SourceLeaf{592}
\begin{center}\small Göttingische gelehrte Anzeigen. Stück 46. Seite 449..456. 1815 März 23.\\[3pt]\rule{.72\linewidth}{.3pt}\end{center}

I. \emph{Neue Methode, beobachtete Azimuthe zu reduciren. Von} J. \textsc{Soldner}. München. 14 S. in Quart. (Aus den Denkschriften der Münchner Academie für 1813 besonders gedruckt.)

II. \emph{Beiträge zur Berechnung beobachteter Azimuthe von} \textsc{Anton von Stefenelli}. Ebendaselbst. Gedruckt mit Stornoschen Schriften. 30 S. in Quart.

Wegen der ungleichförmigen Aenderung der Azimuthe ist es nicht erlaubt, das Mittel aus einer Reihe von Azimuthalunterschieden zwischen einem irdischen Gegenstande und einem Himmelskörper demjenigen Azimuthalunterschiede gleich zu setzen, welcher dem Mittel aller Beobachtungszeiten entspricht. Ist freilich nach jeder einzelnen Beobachtung abgelesen, so wird man auch die Mühe nicht scheuen, jede einzeln zu berechnen, um eben so viele besondere Resultate zu haben. Allein in der Regel pflegt man nur am Ende einer Beobachtungsreihe abzulesen, um lieber die Beobachtungen desto zahlreicher zu machen; die einzelne Berechnung einer grossen Menge von Beobachtungen würde eine abschreckend weitläuftige Arbeit werden, und die Verwechslung des Mittels aller beobachteten Azimuthalunterschiede mit dem Azimuthalunterschiede für das Mittel aller Beobachtungszeiten, würde Fehler hervorbringen, welche bei der Feinheit der Beobachtungen mit Theodolithen, wie sie von \textsc{Reichenbach} gegenwärtig geliefert werden, nicht übersehen werden dürfen. Zum Theil wegen dieses Umstandes schlug daher ein berühmter praktischer Astronom vor, diese Beobachtungen nur um die Zeit der Culmination anzustellen, wo allerdings die Aenderung des Azimuths als gleichförmig betrachtet werden darf, und wo man noch den Vortheil hat, dass Fehler in der Polhöhe oder in der Declination des Gestirns nur einen ganz unmerklichen Einfluss auf das Resultat haben. Indessen treten bei diesem Verfahren wieder andere nachtheilige Umstände ein. Je grösser die Höhe des beobachteten Gestirns ist, desto grösser wird der Einfluss kleiner Fehler bei der Berichtigung des Instruments, die an sich schon sehr delicat ist, und leicht wieder etwas gestört wird, wenn man nicht dem Instrumente eine sehr feste Aufstellung geben kann; ein Fehler in der Zeitbestimmung wirkt hier nachtheiliger auf das Azimuth, als bei kleinen Höhen; bei Beobachtung der Sonne ist man, in so fern das Fernrohr nur bis zu $40^{\circ}$ geneigt werden kann, in unsern oder südlichern Breiten auf die Wintermonate beschränkt; Fixsterne erster Grösse lassen sich zwar auch bei Tage mit jenen Werkzeugen beobachten, allein wenn man die Zeitbestimmung von correspondirenden Sonnenhöhen hernehmen muss, so kommt dann auch noch der etwaige Fehler der Sonnentafeln in Betracht, ist man hingegen in Besitz eines guten Passageninstruments, so wird es in den meisten Fällen der Azimuthalbeobachtungen gar nicht bedürfen, da man mit Hülfe dieses Instruments leicht und mit grosser Schärfe ein künstliches Object in die Mittagsfläche bringen kann. Aus diesen Gründen würden wir den Azimuthalbestimmungen ausser der Culmination und in mässigen Höhen den Vorzug geben. In dem Falle, wo bei der Polhöhe oder der Declination des Gestirns noch eine kleine Ungewissheit vorhanden ist, kann man diese leicht unschädlich machen, wenn man nur Beobachtungen auf beiden Seiten und in beinahe gleichen Entfernungen vom Meridian verbindet; und eben so wenig kann die vorhin erwähnte Ungleichförmigkeit der Azimuthaländerungen hier in Betracht kommen, da sich dieselbe leicht und mit grösster Schärfe in Rechnung bringen lässt. In der That ist klar, dass das Azimuth, welches dem Stundenwinkel $t+\theta$ entspricht, sich nach \textsc{Taylor}'s Lehrsatze in eine Reihe $A+B\theta+C\theta\theta+D\theta^3+\text{etc.}$ entwickeln lässt, von welcher, in so fern $\theta$ nur innerhalb enger Grenzen liegt, wenige Glieder hinreichend sein werden. Sind also zusammen $n$ Beobachtungen gemacht, denen die Stundenwinkel $t+\theta$, $t+\theta'$, $t+\theta''$, $t+\theta'''$ u. s. w. entsprechen, so wird of\SourceLeaf{593}fenbar das Mittel aller Azimuthe
\[
\begin{aligned}&= A + \frac{\theta+\theta'+\theta''+\theta'''+\text{etc.}}{n}\cdot B\\[3pt]
&+ \frac{\theta\theta+\theta'\theta'+\theta''\theta''+\theta'''\theta'''+\text{etc.}}{n}\cdot C\\[3pt]
&+ \text{etc.}\end{aligned}
\]
Da es willkürlich ist, welchen Werth von $t$ man hier zum Grunde legen will, so wird es am einfachsten sein, denjenigen zu wählen, welcher dem Mittel aller Beobachtungszeiten entspricht, wodurch $\theta+\theta'+\theta''+\theta'''+\text{etc.}=0$ wird, und folglich das Mittel aller Azimuthe die Form enthält.
\[
A + \frac{1}{n}\,C\,\Sigma\theta\theta + \frac{1}{n}\,D\,\Sigma\theta^3
\]
denn weiter als bis zu den Gliedern der dritten Ordnung wird man in der Ausübung niemals zu gehen brauchen, ja bei weitem in den meisten Fällen wird es an der Correction der zweiten Ordnung $\frac{1}{n}\,C\,\Sigma\theta\theta$ genug sein, statt welcher man sich offenbar auch erlauben darf zu nehmen
\[
2\,C \times \frac{\Sigma\,2\sin\frac{1}{2}\theta^2}{n}
\]
indem der Fehler nur von der vierten Ordnung wird. Man erhält dadurch den Vortheil, die bekannte Hülfstafel zur Reduction von Circummeridianhöhen benutzen zu können, welche sofort $2\sin\frac{1}{2}\theta^2$ in Secunden ausgedrückt gibt. Dass man bei dem dritten Gliede, so wie $D$ aus der Entwickelung nach \textsc{Taylors} Lehrsatz folgt, noch den Factor $\frac{15^3}{206265^2}$ oder $\frac{900^3}{206265^2}$ oder $\frac{9000^3}{206265^2}$ hinzufügen muss, wenn man $\theta$ in Zeit-Secunden, oder Zeit-Minuten, oder Zehnern von Zeit-Minuten ausdrückt, ist von selbst klar.

Das hier kürzlich beschriebene Verfahren ist es, was Herr \textsc{Soldner} in der vorliegenden Abhandlung vorträgt. Es ist so einfach und liegt so nahe, dass man sich wundern muss, dass es mehrern praktischen Astronomen bei derselben oder bei ganz ähnlichen Veranlassungen entgangen ist. Eben deswegen aber verdient Hr. \textsc{Soldner} für die Bekanntmachung desselben den Dank der praktischen Astronomen, um so mehr, da die Entwickelungen auf eine geschickte und elegante Art durchgeführt, und die Endresultate, d. i. die Ausdrücke für die Coëfficienten $B$, $C$, $D$ (denn so weit hat Hr. \textsc{Soldner} die Entwicklung getrieben), in eine ganz geschmeidige Form gebracht sind.

\begingroup\rightskip=0pt plus 3em
Recensent, welcher bei häufigen mit einem \textsc{Reichenbach}'schn
 Theodolithen seit ein Paar Jahren gemachten Azimuthalbeobachtungen eben dieses Verfahren angewandt hat, ohne jedoch das in der Ausübung meistens überflüssige Glied der dritten Ordnung mit in Betracht zu ziehen, hofft, dass es den astronomischen Lesern nicht unlieb sein wird, hier noch einige Bemerkungen über diesen Gegenstand zu finden. Er bediente sich dabei der Formel
\[
2\,C = -\,\frac{\cos\varphi\cdot \operatorname{tang}\,\delta\cdot \sin A^2}{\sin t}
\left(1 + \frac{2\cos\varphi\cdot \sin 2A}{\sin t\cdot \sin 2\delta}\right)
\]
wo $\varphi$, $\delta$, $t$, $A$ Polhöhe, Declination, Stnndenwinkel
 und Azimuth bedeuten, und die man auch so darstellen kann:
\[
a = \frac{\cos\varphi\cdot \operatorname{tang}\,\delta\cdot \sin A^2}{\sin t}, \qquad
2\,C = -\,a\left(1 + \frac{2a}{\operatorname{tang}\,A\cdot \sin\delta^2}\right)
\]
Diese Formel ist eigentlich identisch mit einer auch von Hrn. \textsc{Soldner} S. 12 angeführten, welche dort auch nur für den Fall empfohlen wird, wo man das Glied der dritten Ordnung nicht mit in Rechnung nehmen will. Nur ist dort nicht bemerkt, dass es \emph{unter dieser Voraussetzung} bequemer ist, zur Berechnung des Azimuths selbst, anstatt der \textsc{Neper}'schen Formeln folgende (an sich bekannte) zu gebrauchen:
\SourceLeaf{594}
\[
\begin{gathered}
\frac{\operatorname{tang}\delta}{\cos t}=\operatorname{tang}\psi\\[5pt]
\frac{\cos\psi\cdot\operatorname{tang}t}{\sin(\varphi-\psi)}=\operatorname{tang}A
\end{gathered}
\]
Soll hingegen das Glied der dritten Ordnung mit beachtet werden, so wird man sich für $C$ und $D$ der \textsc{Soldner}'schen Formeln S. 9 bedienen, an deren Form man auch noch eine kleine unten zu berührende Abänderung anbringen kann.\par\endgroup

So oft man sich mit der Correction der zweiten Ordnung begnügen kann, lässt sich selbst die Berechnung des Coëfficienten $C$ ganz umgehen, vermittelst eines kleinen Kunstgriffs, von welchem man auch bei mancherlei andern Gelegenheiten mit Vortheil Gebrauch machen kann, und den wir hier noch anführen wollen. Es sei, wie vorhin, $t$ der Stundenwinkel für das Mittel aller Beobachtungszeiten, oder $\theta+\theta'+\theta''+\theta'''+\text{etc.}=0$, und
\[
\tau=\BareRoot{\frac{\theta\theta+\theta'\theta'+\theta''\theta''+\theta'''\theta'''+\text{etc.}}{n}}
\]
Das Mittel aller $n$ Azimuthe ist dann gleich zu setzen dem Mittel aus den zwei Azimuthen, welche den Stundenwinkeln $t-\tau$ und $t+\tau$ entsprechen, und die man ungefähr eben so leicht direct berechnen kann, als das eine Azimuth $A$ und den Coëfficienten $C$. Zur leichtern Bestimmung von $\tau$ darf man sich ohne Bedenken der schon oben erwähnten Reductionstafel bedienen, in welcher zu $\tau$ in Zeit ausgedrückt das Mittel der einzelnen $n$ Grössen gehören wird, die in derselben den Argumenten $\theta$, $\theta'$, $\theta''$, $\theta'''$ u. s. w. in Zeit entsprechen. Dies Verfahren empfiehlt sich auch dadurch, dass man dabei, ohne weitere Vergrösserung der Arbeit, auf allerlei Nebenumstände Rücksicht nehmen kann, z. B. bei gegenwärtiger Aufgabe, wenn von der Sonne die Rede ist, auf deren Declinationsänderung. Unmittelbar ist übrigens dieser Kunstgriff nur dann anwendbar, wenn die Stundenwinkel bekannt sind, und die Azimuthe gesucht werden; bei der Anwendung beobachteter Azimuthalunterschiede zur Zeitbestimmung müssten noch andere Hülfsmittel beigefügt werden, zu deren Auseinandersetzung hier nicht der Ort ist.

Betrachtet man die von Hrn. \textsc{Soldner} abgehandelte Untersuchung aus einem blos mathematischen Gesichtspunkte, so gibt sie eigentlich die Auflösung folgender Aufgabe: Wenn in einem sphärischen Dreiecke zwei Seiten unveränderlich sind, die endlichen Aenderungen eines anliegenden Winkels durch eine nach den Potenzen der endlichen Aenderungen des eingeschlossenen Winkels fortlaufende Reihe auszudrücken. Das Gegenstück zu dieser Aufgabe wäre die, unter denselben Bedingungen die endlichen Aenderungen der dritten Seite auf eine ähnliche Art darzustellen, eine Aufgabe, deren Auflösung gleichfalls zur scharfen Reduction von Beobachtungen mit Vervielfältigungswerkzeugen bei mehrern Gelegenheiten nothwendig ist. Mit diesen beiden Aufgaben und verschiedenen Anwendungen beschäftigt sich die Schrift Nro. II. Ihrem Verfasser war die Abhandlung von Hrn. \textsc{Soldner} schon bekannt, wie man aus verschiedenen zum Theil sehr unpassenden Aeusserungen über dieselbe sieht. Dahin gehört S. 29 die seltsame Behauptung, dass das \textsc{Soldner}'sche Verfahren nichts weiter sei, als die von \textsc{Svanberg} angewandte Methode, ein Urtheil, dessen Grundlosigkeit jedem, der beide Methoden vergleicht, zu einleuchtend ist, als dass wir dabei länger verweilen sollten. In Rücksicht der Ausführung können wir der \textsc{Stefenelli}'schen Schrift nicht dasselbe Lob beilegen wie der \textsc{Soldner}'schen. Die beiden Reihen, wodurch Herr \textsc{Stefenelli} die beiden Aufgaben aufzulösen glaubt, und die, wenn wir uns recht erinnern, schon vor einiger Zeit in einem gelehrten Blatt unter dem unpassenden Titel von neuen Entdeckungen angekündigt wurden, sind beide schon in den Gliedern der dritten Ordnung \emph{unrichtig}. Es ist keine Entschuldigung für den Verf., dass diese Glieder in der Ausübung meistens unbedeutend sind; er mochte sich auf die Glieder der zweiten Ordnung beschränken, aber wenn er es einmal unternahm, darüber \SourceLeaf{595} hinaus zu gehen, so fordert man richtige Resultate, um so mehr, da die eine Aufgabe schon von Hrn.
\textsc{Soldner} bis zur dritten Ordnung vollständig aufgelöst war, und überhaupt das ganze Verdienst
dieser an sich gar nicht schweren Entwicklungen lediglich in ihrer Richtigkeit und Geschmeidigkeit
besteht. Das einzige, was wir zum Lobe der \textsc{Stefenelli}'schen Schrift anführen können, ist,
dass dort das noch richtige Glied der zweiten Ordnung (bei der ersten Aufgabe, welche bei Hrn.
\textsc{Stefenelli} die zweite ist) eine etwas bequemere Gestalt hat, als bei Hrn. \textsc{Soldner};
Hr. \textsc{Stefenelli} muss aber nicht bemerkt haben, dass sich der \textsc{Soldner}'schen Reihe
durch die leichten und auf bekannten Gleichungen beruhenden Substitutionen
\[
\begin{gathered}
\sin\varphi+\sin\delta=\frac{2\cos\frac{1}{2}z^2\cdot\sin2\gamma}{\sin t}\\[5pt]
\sin\varphi-\sin\delta=\frac{2\sin\frac{1}{2}z^2\cdot\sin2\beta}{\sin t}
\end{gathered}
\]
sogleich dieselbe Form geben lässt: denn hätte er dies bemerkt, so würde er sich bei der
Vergleichung des Gliedes der dritten Ordnung von der Unrichtigkeit des seinigen haben belehren
können. Man begreift kaum, wie er übersehen konnte, dass die Verwechslung von Bögen mit ihren
Sinus und Tangenten S.\,3 nothwendig schon Fehler der dritten Ordnung hervorbringen musste.

\begingroup\rightskip=0pt plus 3em
Einer Schlussanmerkung zu Folge hat Hr. \textsc{Stefenelli} einen Theil seiner Schrift der
mathematisch-physischen Classe der Münchner Academie vorgelegt, wie es scheint, selbst mit der
Erwartung, dieselbe in die Denkschriften der Academie aufgenommen zu sehen. Dies ist nun freilich
nicht geschehen: ob die Academie ein Urtheil darüber abgegeben hat, wird nicht gesagt. Wir würden
bei Anzeige derselben weniger strenge gewesen sein, oder sie lieber ganz ignorirt haben, da sie so
viel wir wissen des Verfassers erster Versuch zu sein scheint, wenn nicht sein unüberlegter Ausfall
auf die \textsc{Soldner}'sche Abhandlung eine ernstliche Rüge verdiente. Um des Verfassers selbst
willen ist es zu bedauern, dass er nicht veranlasst wurde seinem Versuche erst mehr Reife zu geben,
ehe er damit öffentlich auftrat.\par\endgroup
